\section{Introducción} \label{sec:1_introduccion}

La imagen digital ocupa hoy un lugar muy importante en prácticamente cualquier ámbito
tecnológico. Desde la fotografía computacional hasta la visión artificial, pasando
por la instrumentación científica, la monitorización industrial o los sistemas de
inspección, las imágenes se han convertido en una de las principales formas de
representar, interpretar y comunicar información sobre el entorno. Sin embargo,
la aparente inmediatez con la que una imagen puede observarse en pantalla suele
ocultar una cadena de decisiones y transformaciones mucho más compleja de lo que
parece a simple vista.

Cuando una imagen procede de un flujo de captura cruda, esa complejidad se hace
todavía más evidente. A diferencia de una imagen ya procesada, el dato
crudo conserva una representación más próxima a la adquisición original del
sensor, lo que le confiere un gran valor desde el punto de vista informativo, pero
también introduce nuevas exigencias de tratamiento. En este contexto, cuestiones
como la legibilidad visual, la percepción del detalle o el aprovechamiento del
rango tonal dejan de depender únicamente de la escena capturada y pasan a estar
condicionadas por la forma en que dicha información es interpretada, revelada y
ajustada.

Entre todos los aspectos que influyen en esa interpretación, el contraste ocupa
una posición especialmente relevante. No se trata solamente de una cualidad
estética, sino de un factor determinante en la capacidad de una imagen para hacer
visible su contenido, separar estructuras, destacar regiones de interés o facilitar
su análisis posterior. En los flujos de trabajo basados en imagen cruda, 
esta cuestión adquiere una dimensión particular, ya que la mejora del contraste 
no puede entenderse únicamente como una operación aislada, sino como parte de una secuencia 
de transformaciones estrechamente relacionadas entre sí.

Al mismo tiempo, la necesidad de intervenir sobre la imagen de manera inmediata y
significativa es cada vez mayor. En numerosos escenarios actuales no basta con
obtener una representación visual correcta. Es necesario explorar variantes,
comparar resultados y reajustar parámetros con fluidez, manteniendo una relación
directa entre la intervención del usuario y la respuesta visual del sistema. Esta
dimensión interactiva desplaza el foco desde el procesado entendido como una fase
cerrada hacia una concepción más dinámica, donde la experimentación, la
reproducibilidad y la capacidad de iteración adquieren mucha importancia.

En esta convergencia entre captura cruda, tratamiento de imágenes e interacción 
en tiempo real se sitúa el presente Trabajo Fin de Máster. El proyecto se enmarca 
en el estudio de herramientas y enfoques capaces de acercar el procesamiento digital 
de imágenes a un entorno más inmediato, explorable y controlable, en el que la imagen 
no se contempla únicamente como un resultado final, sino también como un objeto de 
trabajo sobre el que observar, comparar y decidir. Desde esta perspectiva, 
se desarrolla un recorrido que parte del fundamento teórico del problema y culmina 
en la materialización de una solución orientada a hacer del tratamiento de imágenes 
crudas un proceso más comprensible, más flexible y más cercano a las necesidades 
reales de análisis y ajuste visual.
